%==============================================================================
%  CHAPTER 1 -- INTRODUCTION
%
%  Job of this chapter: convince the reader that the question you answered was
%  worth asking, and that nobody had answered it yet.
%  Typically 8-12 pages. Write it second-to-last, once you know what you found.
%==============================================================================
\chapter{Introduction}
\label{ch:introduction}

%------------------------------------------------------------------------------
\section{Background and motivation}
\label{sec:background}
%  Start wide, then narrow. The engineering context, and why anyone should
%  care: what fails, what it costs, who is affected.
%  End the section with the reason this work was undertaken.
%  Write for a competent engineer who does NOT work in your specific field.

%------------------------------------------------------------------------------
\section{State of the art}
\label{sec:state-of-the-art}
%  What has already been done, by whom, and what each approach achieved.
%  This is a CRITICAL review, not a list. Every paragraph should end with what
%  the cited work could not do -- those limitations justify your own work.
%
%  Cite with \cite{key} or \cite{key1,key2}.
%
%  Weak:   "Rossi et al. studied vibration. Bianchi et al. studied fatigue."
%  Strong: "Rossi et al. \cite{rossi2020} resolved the field up to 2 kHz, but
%           the method needs markers that occlude the surface under test."

%------------------------------------------------------------------------------
\section{Open research question}
\label{sec:research-question}
%  State the gap in one or two sentences, then the question itself.
%  Be specific enough that it can be answered yes/no, or with a number.
%
%  Weak:   "This thesis studies thermoelastic measurements."
%  Strong: "Can stress mode shapes be identified on a rotating component
%           without fiducial markers, and up to what rotational speed?"
%
%  Then list the objectives that follow from it -- three to five is normal.

%------------------------------------------------------------------------------
\section{Structure of the thesis}
\label{sec:structure}
%  One short paragraph per chapter, in order, saying what it contains.
%  Keep it factual. This is the map the examiner uses to navigate.
